\section{Tools, Skills ו-Agent Skills}

סוכן \eng{AI} מבדיל בין \textbf{יכולת ביצוע} (\eng{Tool}) לבין \textbf{מעטפת
ידע והנחיה} (\eng{Skill}). הבנה נכונה של ההבחנה מאפשרת לבנות מערכות מודולריות
שניתן לתחזק ולהרחיב \cite{yao2023react,openai2024agents}.

\subsection{Tools --- כלי ביצוע}

\eng{Tools} הם התוכנות והפונקציות שהסוכן מפעיל: חיפוש באינטרנט, הרצת
\eng{Python}, קריאת קובץ, קומפילציית \LaTeX{}, שליחת \eng{HTTP}. ה-\eng{LLM}
לא ``מריץ'' קוד בעצמו --- הוא מחליט \emph{איזה} כלי להפעיל ועם אילו פרמטרים,
והסביבה מבצעת בפועל.

\subsection{Skills --- קבצי הנחיה}

\eng{Skill} הוא קובץ (לרוב \eng{Markdown}) המגדיר תפקיד, כללים ודוגמאות.
בתחילת כל \eng{Skill} מופיע \eng{description} שמסביר מתי להשתמש בו. הסוכן
סורק את רשימת ה-\eng{Skills}, בוחר את הרלוונטיים לפרומפט וטוען את תוכנם
לחלון ההקשר --- בדומה לבחירת מומחה מספרייה.

\begin{notebox}
במטלה 02 בקורס נבנו \eng{Skills} נפרדים לכל סוכן תרגום
(\texttt{skills/en\_to\_fr.md} וכו'). אותו עיקרון חל על ייצור \eng{PDF},
\eng{QA} וחקירת סייבר.
\end{notebox}

\subsection{Agent Skills --- סקילים ברמת הסוכן}

\eng{Agent Skills} (כמו ב-\eng{Cursor} או \eng{Claude}) הם מפרטים שניתן
לשתף בין פרויקטים: הוראות מערכת, תבניות פרומפט וכללי בטיחות. הם מהווים
\eng{Memory} פרוצדורלי --- ידע שתמיד זמין לסוכן בלי לאחסן אותו בכל שיחה.

\subsection{תהליך בחירת Skill}

לפי ההרצאה, כשמשתמש שואל שאלה:
\begin{enumerate}[rightmargin=2em]
  \item ה-\eng{LLM} סורק תיאורי \eng{Skills} זמינים.
  \item נבחרים \eng{Skills} עם התאמה סמנטית (למשל ``חיפוש באינטרנט'').
  \item תוכן ה-\eng{Skill} נטען ל-\eng{Context Window}.
  \item הסוכן מפעיל \eng{Tools} שה-\eng{Skill} מגדיר או מרמז עליהם.
\end{enumerate}

\subsection{דוגמה: Skill לקומפילציית LaTeX}

\begin{lstlisting}[caption={קטע מתוך Skill לייצור PDF}]
# LaTeX Compile Skill
When the user requests a PDF document:
1. Use agentic-ai-template.cls (do not inline styles in main.tex)
2. Compile with: lualatex -> bibtex -> lualatex x2
3. Fix recurring RTL issues in the .cls file, not per-document
4. Verify TOC and bibliography before delivery
\end{lstlisting}

\subsection{Skills מול Tools --- סיכום}

\agenttable{
\caption{הבחנה בין \eng{Tool} ל-\eng{Skill}}
\label{tab:tools-skills}
\begin{tabular}{@{}p{2.5cm}p{4.5cm}p{4.5cm}@{}}
\toprule
 & \textbf{Tool} & \textbf{Skill} \\
\midrule
מהות & קוד רץ & הנחיות טקסט \\
מיקום & \eng{src/}, \eng{MCP} & \texttt{skills/*.md} \\
שינוי & מפתח תוכנה & מומחה תחום \\
דוגמה & \eng{vector\_compare.py} & \eng{QA-Table.md} \\
\bottomrule
\end{tabular}
}

\begin{insightbox}
\eng{ReAct} \cite{yao2023react} משלב חשיבה (\eng{Reasoning}) ופעולה
(\eng{Acting}) --- המודל כותב מחשבות, בוחר כלי, מקבל תוצאה וחוזר חלילה.
זהו הבסיס לרב-סוכנים מודרניים.
\end{insightbox}
